% ============================================================
% Plantilla: Manuscrito para revistas de Elsevier (clase oficial elsarticle)
% Opciones útiles: preprint (una columna, para enviar), review (doble espacio),
% 1p / 3p / 5p (formatos de publicación). Compila con pdflatex + bibtex.
% ============================================================
\documentclass[preprint,12pt]{elsarticle}

\usepackage[utf8]{inputenc}
\usepackage[T1]{fontenc}
\usepackage[spanish,es-tabla]{babel}
\usepackage{amsmath,amssymb}
\usepackage{booktabs}

\journal{Nombre de la Revista}

\begin{document}

\begin{frontmatter}

\title{Título del manuscrito}

\author[uni]{Primer Autor\corref{cor1}}
\ead{autor@ejemplo.com}
\author[uni,inst]{Segundo Autor}

\cortext[cor1]{Autor de correspondencia.}

\affiliation[uni]{organization={Facultad de Ingeniería, Universidad},
  city={Ciudad}, country={País}}
\affiliation[inst]{organization={Instituto de Investigación},
  city={Ciudad}, country={País}}

\begin{abstract}
Resume el objetivo, la metodología, los resultados y las conclusiones del trabajo.
\end{abstract}

\begin{highlights}
\item Primera contribución principal del trabajo.
\item Segunda contribución principal.
\item Tercera contribución principal.
\end{highlights}

\begin{keyword}
palabra clave uno \sep palabra clave dos \sep palabra clave tres
\end{keyword}

\end{frontmatter}

\section{Introducción}
Presenta el problema y la contribución del manuscrito \cite{lamport1994}.

\section{Materiales y métodos}
Describe el procedimiento con detalle suficiente para reproducirlo \cite{knuth1984}.

\section{Resultados}
\begin{table}[htbp]
  \centering
  \caption{Resultados principales.}
  \begin{tabular}{lcc}
    \toprule
    Muestra & Valor A & Valor B \\
    \midrule
    M1 & 12.4 & 0.81 \\
    M2 & 15.1 & 0.87 \\
    \bottomrule
  \end{tabular}
\end{table}

\section{Conclusiones}
Resume los hallazgos \cite{goossens2004}.

\bibliographystyle{elsarticle-num}
\bibliography{referencias}

\end{document}
